The synthesis of \cite{SeventeenthPaper2026} traced the stalemate on whether a multi-tape
Turing machine can transpose a $d\times d$ matrix of $\clog{d}$-bit entries in $O(L)$ steps,
where $L = d^2\clog{d}$ (question $(\T)$), to a structural wall: trace reductions on the
machine side produce bounded-fan-in cut cones ($\Delta = O_{k,\eps}(1)$) while discarding
the finite transition table that generates them, whereas shared-decoder models retain a
shared $M$-class rule while discarding bounded read arity. Because the two invariants were
analyzed separately, arbitrary-set shared decoders blew up to rate
$1 - \frac{\log_{\abs{\Sigma}} d}{d^2} \to 1$ at unrestricted arity, two-stage linear
composition relied on Hales--Jewett arithmetic progressions and collapsed at
$M_1 \approx \log\log d$, and cyclic column-shift symmetry could not distinguish matrix
transposition from Stone's $O(L)$-time perfect shuffle.

This paper reunites bounded read arity with decoder sharing and establishes six main
kernel-verified results in Lean~4. First (\Cref{thm:set-capacity}), for every alphabet
$\abs{\Sigma} \ge 2$ (with no dependence on $\abs{\Sigma}$), any subset
$C \subseteq \Sigma^{d\times d}$ transposed outside $\mathcal{E}$ by $M$ shared deterministic
column decoders of cell read arity $\Delta \ge 1$ satisfies
\[
  \log_{\abs{\Sigma}}\abs{C} \;\le\; d^2 \;-\;
  \bigl(d^2 - \abs{\mathcal{E}} - dM\Delta\bigr)\frac{\Delta^\Delta}{(\Delta+1)^{\Delta+1}},
\]
giving the first unconditional strictly-below-$1$ capacity bound on arbitrary nonlinear sets
at every constant arity $\Delta \ge 3$ (where \cite[Problem~9.4(b)]{SeventeenthPaper2026} remained
open at rate $1$, giving e.g.\ $\le \frac{229}{256} + o(1)$ at $\Delta = 3$), the first Lean-verified
bound below $1$ at $\Delta \ge 2$ (complementing the written $M = 1, \Delta \le 2, \mathcal{E}=\emptyset$
order bound of \cite[Proposition~6.22(b)]{SeventeenthPaper2026}), and a positive rate deficit across
all $\Delta \ge 1$ whenever $M\Delta = o(d)$ and $\abs{\mathcal{E}} \le (1-\delta)d^2$. Second (\Cref{thm:2a,thm:2b}), we eliminate
Hales--Jewett Ramsey numbers from two-stage linear shared composition: via affine Hilbert
cubes over any field $F$ we prove the explicit exact-factorization cutoff
$d \le \max\bigl(8 M_1 M_2^2,\; (2M_1)^{2^{M_2}}\bigr)$ (ruling out $K_2 \circ K_1 = \tau$ up to
$M_1 = o(d^{1/2^{M_2}})$ for fixed $M_2$, including on $3$-AP-free Behrend partitions), and over
any field with $(d : F) \ne 0$ we combine shifted sumsets, arbitrary-subset Vandermonde pullbacks,
sparse annihilating polynomials, and iterated K\H{o}v\'ari--S\'os--Tur\'an extraction to prove an
explicit rank deficit $\Omega\!\bigl(b\,d\,(\log d)^{1/(2M_2-1)}\bigr)$ for all constant $M_1, M_2$
($\Omega(b\,d\log d)$ when $M_2 = 1$). Third (\Cref{thm:anti-diagonal,thm:osr-pass}),
anti-diagonal untwisting identifies row-shift-equivariant operators with decoupled tensor
slices---packaging the $9/16$ and $5/8$ random-mask certificate floors and realizing the
single-shift $5/8$ floor via generic tensor subspaces---while Operator Schmidt Rank
across the row/column tensor cut separates the transpose ($\OSR(\tau) = d^2$) from Stone's
shuffle ($\OSR(S) = 4$ for $d = 2^m \ge 2$), yielding a $t \ge \log_2 d$ pass lower bound on alternating
shuffle/row/column pipelines (\cite[Problems~9.1 and~9.5]{SeventeenthPaper2026}). Fourth
(\Cref{thm:three-term-unified,thm:multiscale-interval,thm:walsh-dichotomy}), we disprove the
$d(d+1)/2$ pairing bound for three-term swap codes asked in
\cite[Remark~5.6]{SeventeenthPaper2026} (attaining dimension $4m(m+1)$ at $d = 3m$, which exceeds
$d(d+1)/2$ for $m \in \{1,2,3,4\}$, and $29 > 28$ at $d = 7$, giving $d \in \{3,4,6,7,9,12\}$ over
every field), resolve the literal multi-scale interval formulation of
\cite[Problem~9.3]{SeventeenthPaper2026} with exact upper bound $(d^2 + d\beta)/2$ where
$\beta = \min\bigl(\max(1, 2\sum_j M_{A,j}),\, \max(1, 2\sum_j M_{B,j})\bigr)$,
and isolate the periodic Rademacher/Walsh boundary dichotomy. Finally
(\Cref{thm:sweep-cut-absorbed,thm:sweep-cut-refined}), we construct native causal-cone
evaluation trees directly from the transition function of \leanname{Transpose.TM} on frozen
cuts (\cite[Problem~9.4]{SeventeenthPaper2026}). Question $(\T)$ itself remains open; we
conclude by isolating the sharp finite and algebraic thresholds to which the series now
reduces $(\T)$.
